% Part: computability
% Chapter: recursive-functions
% Section: bounded-minimization

\documentclass[../../../include/open-logic-section]{subfiles}

\begin{document}

% SEGMENT OLP-0218-S01

\olfileid{cmp}{rec}{bmi}
\olsection{வரம்பிட்ட சிறுமமாக்கல்}

\begin{explain}
ஏதேனும் பண்பு அல்லது உறவு~$P$ ஐ நிறைவுறுத்தும் மிகச்சிறிய எண்ணாக
ஒரு சார்பை வரையறுப்பது பலவேளைகளில் பயனுள்ளதாகும். $P$
தீர்மானிக்கத்தக்கது எனில்,~$P$ ஐ நிறைவுறுத்தும் மிகச்சிறிய எண்ணைக்
காணும் வரை சாத்தியமான அனைத்து எண்களையும் $0$, $1$, $2$, \dots,
எனச் சோதிப்பதன் மூலம் இச்சார்பைக் கணிக்கலாம். இத்தகைய வரம்பற்ற தேடல்
நம்மை முதன்மை மீள்வரையறைச் சார்புகளின் எல்லைக்கு வெளியே கொண்டு
செல்கிறது. எனினும், \emph{சார்பின்றித் தரப்பட்ட ஏதேனும் வரம்பைவிடக்
குறைந்த} மிகச்சிறிய எண்ணில் மட்டும் ஆர்வம் இருந்தால், நாம் முதன்மை
மீள்வரையறைக்குள் இருக்கிறோம். வேறுவிதமாகவும் சற்றுப் பொதுவாகவும்
கூறினால், முதன்மை மீள்வரையறை உறவு~$R(x,z)$ ஒன்று நம்மிடம் உள்ளது
என்று கொள்க. $R(x, z)$ பொருந்தும் மிகச்சிறிய $z < y$ க்கு
$x$ மற்றும்~$y$ ஐக் கொண்டுசெல்லும் சார்பைக் கருதுக. $R(x, 0)$,
$R(x, 1)$, \dots, $R(x, y-1)$ ஆகியவை பொருந்துகின்றனவா என்று
சோதிப்பதன் மூலம் இதையும் கணிக்கலாம். ஆனால் இது ஏன் முதன்மை
மீள்வரையறைச் சார்பாகிறது?
\end{explain}

\begin{prop}
$R(\vec x, z)$ ஒரு முதன்மை மீள்வரையறை உறவு எனில்,
$R(\vec x, z)$ பொருந்தும் மிகச்சிறிய~$z$ ஆனது~$y$ ஐவிடக் குறைந்து
இருந்தால்
அதையும், இல்லையெனில் $y$ ஐயும் திருப்பித்தரும்
$m_R(\vec{x}, y)$ என்ற சார்பும் முதன்மை மீள்வரையறைச் சார்பாகும்.
சார்பு~$m_R$ ஐப் பின்வருமாறு எழுதுவோம்:
\[
\bmin{z < y}{R(\vec{x}, z)},
\]
\end{prop}

\begin{proof}
$R(\vec{x}, z)$ பொருந்துமாறு~$z < 0$ எதுவும் இருக்க முடியாது;
ஏனெனில் $z < 0$ என்பதே எதுவும் இல்லை. எனவே
$m_R(\vec x, 0) = 0$.

வரம்பு $y + 1$ என்ற வடிவில் இருக்கும்போது நமக்கு மூன்று நேர்வுகள்
உள்ளன:
\begin{enumerate}
\item $R(\vec{x}, z)$ பொருந்துமாறு $z < y$ ஒன்று உள்ளது; இந்நேர்வில்
$m_R(\vec{x}, y+1) = m_R(\vec{x}, y)$.
\item அத்தகைய~$z<y$ இல்லை; ஆனால் $R(\vec{x}, y)$ பொருந்துகிறது;
அப்போது $m_R(\vec{x}, y+1) = y$.
\item $R(\vec{x}, z)$ பொருந்துமாறு $z < y+1$ எதுவும் இல்லை;
அப்போது $m_R(\vec{z}, y+1) = y+1$.
\end{enumerate}
எனவே முதன்மை மீள்வரையறையால் $m_R(\vec x, 0)$ ஐப் பின்வருமாறு
வரையறுக்கலாம்:
\begin{align*}
m_R(\vec{x}, 0) & = 0\\
m_R(\vec{x}, y+1) & =
\begin{cases}
m_R(\vec{x}, y) & \text{$m_R(\vec x, y) \neq y$ எனில்}\\
y & \text{$m_R(\vec x, y) = y$ மற்றும் $R(\vec{x}, y)$ எனில்}\\
y+1 & \text{இல்லையெனில்.}
\end{cases}
\end{align*}
$R(\vec x, z)$ பொருந்துமாறு $z<y$ ஒன்று உள்ளது எனில், எனில்
மட்டுமே, $m_R(\vec x, y) \neq y$ என்பதைக் கவனிக்க.
\end{proof}

\begin{prob}
$R(\vec x, z)$ ஒரு முதன்மை மீள்வரையறை உறவு என்று கொள்க.
$R(\vec{x}, z)$ பொருந்தும் மிகச்சிறிய~$z$ ஆனது~$y$ ஐவிடக் குறைந்து
இருந்தால்
அதையும், இல்லையெனில் $0$ ஐயும் திருப்பித்தரும்
$m'_R(\vec{x}, y)$ என்ற சார்பை~$\Char{R}$ இலிருந்து முதன்மை
மீள்வரையறையால் வரையறுக்க.
\end{prob}

\end{document}
